\chapter{HTTPS Credentials}\label{chap: Http Credentials}

When you clone a repository to your PC using the HTTPS protocol (the default when you click the ``Download Code'' button) you will be able to clone the repository to your PC without any problems.

When you make and commit changes, attempting to \code{git push} them back to your repository will ask for your username and password. The push will fail with a message telling you:

\begin{lstlisting}
remote: Invalid username or token.
Password authentication is not supported for Git operations.
fatal: Authentication failed for
'https://github.com/<username>/<repository name>.git/'
\end{lstlisting}

The solution is to either use SSH, See \chapref{chap: Ssh Credentials} or create a Personal Access Token or PAT, as follows:

\begin{itemize}
    \item Login to your account on GitHub.
    \item Click your profile picture on the far right.
    \item Click ``Settings''.
    \item On the left sidebar, click ``Credentials''.
    \item Back to the right, click ``Personal Access Token (Classic)''.
    \item Click ``Generate New Token''
    \item Click ``Generate New Token (Classic)''
    \item Add a note of what it is for. Be meaningful!
    \item Set an expiry date, as appropriate.
    \item Tick ``repo'' \& ``Workflow'' at the very least. You will also want to add some of the other options as appropriate, \emph{read them carefully!} I suggest ``gist'', ``Notifications'' and \emph{maybe} ``delete\_repo'' (oooh! Brave!)
    \item Scroll down and click ``Generate token''.
\end{itemize}

\begin{tips}{Warning}
    The next screen to appear will list the generate PAT for your use. Make a copy of it \emph{right now} because you will never see it again!
\end{tips}

Back on your laptop or PC, you can try the \code{git push} exercise again. This time, when prompted for a password, you must supply the PAT you generated instead.

However, before doing that, do this:

\begin{lstlisting}
git config [--global] credential.helper store
\end{lstlisting}

If you leave ``--global'' out of the command, the credential will be stored for the current repository clone only. Otherwise, it will apply to any and all repositories from the username you generated it under.

Now when you type the password (and where you use the PAT as the password) it will be stored and you won't need to remember it in future.

\begin{tips}{Warning}
    The PAT will be stored on your computer in \emph{plain text} so consider the security implications is you decide to store the PAT.
\end{tips}

